\section{Context of the internship}
\label{sec:context}

% GOAL (2-3 pages). Say where you were, with whom, on what, and how the work
% was organised over time. Keep it factual; the science starts in the next
% section.

%     \begin{itemize}
%         \item Inria: French national research institute for digital science
%               and technology, organised in research centres, each hosting
%               project-teams that are  joint with the university.
%         \item The Inria Centre at the University of Bordeaux:
%               \todoitem{number of teams and staff, main research axes,
%               location on the Talence campus}.
%         \item Local ecosystem relevant to this work: HPC culture of the lab
%     \end{itemize}
\subsection{Inria and the Inria Centre at the University of Bordeaux}

Inria is the French national research institute for digital science and
technology. Its research activities cover a wide range of topics in
computer science, applied mathematics, and digital technologies. Research
is organised into project teams, which bring together researchers,
engineers, PhD students, and other scientific staff around common research
topics.

The internship was carried out at the Inria Centre at the University of
Bordeaux. The centre is located on the Talence campus and brings together
research teams working on different areas of digital science. Inria
currently has 22 project teams associated with the centre, including 20
teams located in Talence and two teams in Pau. The centre brings together
more than 470 researchers, engineers, technical staff, and other members
of the research and innovation community. These teams are developed in
partnership with institutions such as the University of Bordeaux,
Bordeaux INP, CNRS, Inserm, INRAE, and other academic and industrial
partners~\cite{inria_bordeaux}.

The research activities of the centre cover several areas of digital
science. These include high-performance computing and Big Data, machine
learning, and digital health. This diversity reflects the different
research problems addressed within the centre, ranging from fundamental
computer science to applications in science, engineering, and health.
For the work presented in this report, the high-performance computing
research activities are particularly relevant.

Several project teams at the centre work on high-performance and
distributed computing. For example, the STORM, TADaaM, CONCACE, and TOPAL
teams are associated with research topics related to high-performance
computing, systems, optimisation, and the execution of parallel
applications~\cite{inria_hpc_teams}. These teams contribute to a local
research environment in which computing performance, system software, and
large-scale data management are studied using both theoretical approaches
and experiments on real HPC platforms.



%     \item \textbf{The host team.}
%     \begin{itemize}
%         \item Team \textsc{TADaaM}
%         \item Position of this internship inside the team's roadmap: storage
%               performance characterisation, I/O benchmarking, and the
%               development of the IOPS framework.
%         \item Supervision: Luan Teylo.
%         \item People involved in the broader IO500 effort:
%               \todoitem{list the collaborators: people we had the meeting}.
%     \end{itemize}
\subsection{The host team : TADaaM~\cite{inria_tadaam}}

The internship was carried out within the TADaaM project team at the
Inria Centre at the University of Bordeaux. TADaaM works on data
management for high-performance computing systems, with a particular
interest in the interaction between applications, computing resources,
networks, and storage systems. The team studies how data can be managed
efficiently at system scale and how the characteristics of the underlying
infrastructure affect application performance.

The internship is part of the team's activities related to storage
performance characterisation and I/O benchmarking. In particular, the
work contributes to the study of how HPC storage systems can be configured
and evaluated using the IO500 benchmark suite. The internship also makes
use of IOPS, a benchmark orchestration and optimisation framework
developed within the team. The framework provides the tools
needed to automate benchmark executions, explore different parameter
configurations, and collect and analyse their results.

The internship was supervised by Luan Teylo, a researcher
working within the TADaaM team. The work was also connected to a broader
collaboration around IO500 and HPC I/O performance. Discussions involved
Liam Radita, Sara Neuwirth, and Jay Lofstead. These interactions provided
feedback on the benchmark configuration, experimental methodology, and
interpretation of the results.

The specific role of the internship was therefore to investigate the
automated tuning of IO500 configurations. The work focused on the
experimental side of this research activity, from running the benchmark
on the target platform to evaluating automated search strategies and
validating their results.




%     \item \textbf{Scientific and technical environment.}
%     \begin{itemize}
%         \item \textbf{PlaFRIM}: the experimental platform used for all the
%               measurements. \todoitem{describe: number of nodes, the
%               \texttt{bora} partition used here (CPU model, cores per node,
%               memory, interconnect), SLURM as the scheduler}.
%         \item \textbf{Storage}: BeeGFS 
%         \item \textbf{Software stack}: SLURM, OpenMPI
%               \todoitem{4.0.1 was used, confirm}, Python~3.10 or later, IOPS,
%               IO500.
%     \end{itemize}

\subsection{Scientific and technical environment}

The experiments conducted during the internship were carried out on
PlaFRIM, the experimental computing platform used by the Inria research
community in Bordeaux. PlaFRIM provides access to several types of
computing nodes and is designed to support experiments on high-performance
computing systems and new hardware technologies~\cite{plafrim}.

The experiments presented in this report were conducted on the
\texttt{bora} nodes, which consist of 44 nodes, from \texttt{bora001} to
\texttt{bora044}. Each node is equipped with two 18-core Intel Xeon Gold
6240 processors, providing 36 physical CPU cores per node, and 192~GB of
memory. The nodes are connected through a 100~Gbit/s Omni-Path network,
as well as a 10~Gbit/s Ethernet network. For the experiments described in
this report, the BeeGFS parallel file system is accessed through the
100~Gbit/s Omni-Path network~\cite{plafrim}.

The use of a dedicated and homogeneous set of nodes is important for the
experiments carried out in this work. Since the objective is to compare
different IO500 configurations, changes in the measured performance
should primarily result from the parameters being studied rather than
from differences between compute nodes.~\cite{plafrim}.

The experiments were managed using SLURM, the workload manager used on
PlaFRIM. SLURM was used to request computing resources and to execute
benchmark experiments on the selected nodes. Depending on the experiment,
jobs could therefore be allocated with a specified number of nodes and
resources. This was particularly important for the internship because
the number of compute nodes and processes is itself part of the parameter
space explored when tuning IO500.

The storage system used for the experiments is BeeGFS, a parallel file
system designed to provide high-performance access to shared storage.
The benchmark data and results were stored on the BeeGFS file system
provided by PlaFRIM. In addition to BeeGFS, the \texttt{bora} nodes
provide a local 1~TB disk mounted under \texttt{/tmp}. The experiments
described in this report, however, focus on the performance of the
parallel storage system rather than the local disks~\cite{plafrim}.

The software environment consisted of Linux together with the tools
required to build and execute the benchmarks. Python 3.10 was used for
the development and execution of the scripts associated with the
experiments. OpenMPI was used as the MPI implementation required by the
parallel benchmark workloads. The main research tools were IO500, used
to evaluate the I/O performance of the system, and IOPS, used to automate
the execution and exploration of different benchmark configurations.

The development process was split between the PlaFRIM environment and a
local development machine. The source code, scripts, and experimental
configurations were managed using Git and GitLab. This allowed the
development work and experimental results to be kept under version
control while the computationally intensive benchmark executions were
performed on the cluster.



%     \item \textbf{Objectives and specifications.}
%     \begin{itemize}
      %   \item Distinguish what was fixed at the start from what evolved during
            %   the internship, for instance the addition of the specialist
            %   challenge and of the BO diagnostic harness.
%     \end{itemize}

\subsection{Objectives and specifications}


The main objective of the internship was to tune the IO500 benchmark on
the PlaFRIM \texttt{bora} cluster and determine whether automated search
strategies could find good configurations with fewer evaluations than an
exhaustive search.

From the beginning, the main elements of the study were fixed: IO500 was
used as the benchmark, PlaFRIM \texttt{bora} as the target platform, and
IOPS as the framework for running the experiments. Bayesian Optimization
(BO) was selected as the main search strategy, with random search and
exhaustive search used for comparison.

The first stage of the internship focused on the exhaustive search. The
objective was to explore the parameter space and obtain measurements that
could later be used to compare the different search strategies. Each
configuration was also evaluated several times to account for the
variability of HPC measurements.

The study then focused on comparing the different search strategies under
a limited number of evaluations. In particular, the objective was to
determine whether BO could find configurations close to the best
configurations found by exhaustive search while requiring significantly
fewer evaluations.

As the internship progressed, the scope was extended to the development
of IOPS. A plugin for running IO500 experiments was added to IOPS,
together with an example YAML configuration describing the benchmark and
the parameters explored during the study. This made it possible to run
the IO500 experiments directly through IOPS and provided a reproducible
configuration for the experiments.

The IOPS documentation was also updated to describe the new functionality
and its usage.

The search component of IOPS was also extended with an implementation of
the Tree-structured Parzen Estimator (TPE). At this stage, the method had
been implemented but had not yet been evaluated on the IO500 experiments.
Other optimisation methods, such as Iterated Racing (irace) and
Covariance Matrix Adaptation Evolution Strategy (CMA-ES), were also
considered as possible extensions.

These additional tasks emerged as the project evolved and extended the
initial scope of the internship.



%--------------------------------------------------------------------
%     \item \textbf{Organisation of the work.}
%     \begin{itemize}
%         \item Timeline of the internship: bibliography, tooling, exhaustive
%               campaign, offline study, validation campaign, writing.
%               \todoitem{a Gantt chart would work well here}
%         \item Working method: weekly meetings, Git/GitLab repositories,
%               notebooks and reproducible scripts.
%         \item Deliverables produced: the study repository, the IOPS IO500
%               configuration, this report, and the intended IOPS technical
%               report.
%     \end{itemize}
\subsection{Organisation of the work}

The internship took place from 14 April to 16 September. The work was
organised progressively, starting with the study of the existing tools and
the experimental environment, followed by the experimental campaigns and
the development work.

The work was supervised by Luan Teylo through a weekly meeting of
approximately 30 minutes. These meetings were used to review the progress
of the work, discuss experimental results, and define the next tasks.
Notes and follow-up tasks were recorded using the Inria Notes platform.
Mattermost was also used for communication.

The work was managed using Git and GitLab. A personal repository was used
to keep track of the work carried out during the internship, while the
IOPS team repository was used for contributions to the existing IOPS
project. This allowed the development work and experimental scripts to be
versioned throughout the internship.

The internship also included a meeting with Liam Radita, Sara Neuwirth,
and Jay Lofstead as part of the broader collaboration around the IO500
work.

